% alt: A schematic density of the surprisal K = -log p of an interval X: a blue peak at K = S(X) with exponentially decaying tails, dashed vertical lines at S(X) - w and S(X) + w, the right tail beyond S(X) + w shaded orange with its mass bounded by 2 e^{e/2} e^{-u_0 w}, a short double arrow of length 1/u_0 marking the decay scale, and a brace under the axis counting at most e^{S(X)+w} eigenvalues with K <= S(X) + w.
% Section 3 (The marginal tail bound), Lemma lem:tail, eq:tail (part1_arealaw.tex), reused in Section 7 for truncation (part2a_mps_information.tex, proof of Theorem thm:Bshort): K is the surprisal -log p_j drawn with probability p_j, its mean is S(X); the horizontal axis is K, the vertical axis the probability (schematic, no scale); the blue curve is the law of K, the dashed lines are K = S(X) -/+ w, the orange region is the right tail {K > S(X) + w} of mass at most 2 e^{e/2} e^{-u_0 w} (eq:tail, one side), the double arrow is the decay length 1/u_0 of eq:theta depending only on q, R and J/Delta, and the brace under the axis is the count of eigenvalues p of rho_X with -log p <= S(X) + w, at most e^{S(X)+w}, which are the ones kept by the truncation; no tensor is drawn and no leg is contracted.
\begin{tikzpicture}[x=1cm, y=1cm, line width=0.3pt,
    lab/.style={font=\small, inner sep=1pt},
    name/.style={font=\scriptsize, text=black!55, inner sep=1pt},
    brace/.style={decorate, decoration={brace, mirror, amplitude=3pt}},
    cut/.style={dashed, color=orange!80!black, line width=0.5pt}]
  % geometry: the mean S(X) sits at x = 3.4, the half-width w = 1.8, the decay
  % length 1/u_0 = 1 (the curve is 2 sech(K - S(X)), with tails e^{-|K-S(X)|})
  \def\S{3.4}\def\w{1.8}\def\decay{1}\def\xmax{9.6}\def\xmin{-0.4}
  % the right tail beyond S(X)+w, shaded
  \fill[orange!30] plot[domain=5.2:\xmax, samples=60, smooth]
    (\x, {2/cosh(\x-\S)}) -- (\xmax, 0) -- (5.2, 0) -- cycle;
  % hairline range frame
  \draw[black!55] (\xmin, 0) -- (\xmax, 0);
  \draw[black!55] (\xmin, 0) -- (\xmin, 2.2);
  \node[lab, anchor=north east] at (\xmax, -0.05) {surprisal \(K=-\log p\)};
  \node[lab, anchor=south, rotate=90] at (\xmin-0.18, 1.1) {probability};
  % the law of the surprisal
  \draw[blue!65!black, line width=0.6pt] plot[domain=\xmin:\xmax, samples=120, smooth]
    (\x, {2/cosh(\x-\S)});
  \node[name, anchor=south west] at (\S+0.12, 2.0) {mean \(S(X)\)};
  % ticks and the two dashed cuts
  \draw[black!55] (\S, 0.06) -- (\S, -0.06);
  \node[lab, anchor=north] at (\S, -0.08) {\(S(X)\)};
  \draw[cut] (\S-\w, 0) -- (\S-\w, 1.3);
  \draw[cut] (\S+\w, 0) -- (\S+\w, 1.3);
  \node[lab, anchor=north, text=orange!80!black] at (\S-\w, -0.08) {\(S(X)-w\)};
  \node[lab, anchor=north, text=orange!80!black] at (\S+\w, -0.08) {\(S(X)+w\)};
  % the decay length, above the tail
  \draw[black!55, {Bar}-{Bar}, line width=0.4pt] (\S+\w, 0.85) -- (\S+\w+\decay, 0.85);
  \node[lab, anchor=south] at (\S+\w+0.5*\decay, 0.88) {\(1/u_0\)};
  \node[name, anchor=west] at (\S+\w+\decay+0.1, 0.85) {depends only on \(q\), \(R\), \(J/\Delta\)};
  % the tail mass
  \node[lab, anchor=west, text=orange!80!black] at (\S+\w+1.7, 0.42)
    {mass \(\le 2\mathrm{e}^{\mathrm{e}/2}\mathrm{e}^{-u_0w}\)};
  \draw[orange!80!black, -{Stealth[length=3pt]}] (\S+\w+1.65, 0.38) -- (\S+\w+1.1, 0.12);
  % the count under the axis, used for truncation
  \draw[brace, black!55] (\xmin, -0.42) -- (\S+\w, -0.42)
    node[lab, midway, below=4pt] {at most \(\mathrm{e}^{S(X)+w}\) eigenvalues with \(K\le S(X)+w\)}
    node[name, midway, below=16pt] {kept by the truncation};
\end{tikzpicture}
